\thispagestyle{plain}
\begin{center}
    {\fontsize{14}{18}\selectfont\textbf{\underline{ABSTRACT}}}\\[0.8cm]
\end{center}

{\fontsize{12}{18}\selectfont
Modern encrypted voice-over-IP (VoIP) communications rely heavily on classical key exchange mechanisms such as Diffie-Hellman and elliptic-curve cryptography, which are fundamentally vulnerable to large-scale quantum computers running Shor's algorithm. While post-quantum cryptography (PQC) provides mathematical resilience against ``harvest now, decrypt later'' eavesdropping campaigns, its adoption in real-time interactive voice transport introduces a severe structural networking bottleneck: cryptographic key material is exceptionally large. Standardized lattice-based schemes, such as ML-KEM-768 for key encapsulation and ML-DSA-65 for digital signatures, require a combined handshake payload of 10,842 bytes—more than seven times the typical 1,500-byte Internet Maximum Transmission Unit (MTU). 

When transmitted naively across datagram protocols such as UDP, this payload forces the operating system kernel to slice the handshake into eight separate IPv4 fragments. In lossy wireless environments or networks traversing stateless firewalls and network address translation (NAT) middleboxes, losing even a single fragment invalidates the entire reassembly buffer, causing the overall handshake delivery probability to collapse along a steep $(1-p)^8$ exponential curve. Consequently, at an ordinary packet loss rate of 20\%, naive PQC-over-UDP calls establish successfully only 19 times out of 100 attempts.

To resolve this conflict between cryptographic security and real-time network survivability, this thesis introduces \textbf{HPQV} (Hybrid Post-Quantum Voice), a decoupled dual-channel transport protocol. HPQV routes the oversized, mutually authenticated post-quantum handshake, call signaling, and dynamic key rotation ratchets over a reliable, byte-stream TCP control channel, while streaming low-latency Opus audio frames over an independent UDP data channel. Every voice datagram is encapsulated with an 11-byte authenticated header and encrypted via ChaCha20-Poly1305 AEAD, with per-direction nonce prefix derivation to prevent catastrophic nonce reuse, and is bounded by an explicit 1,200-byte application-level guard with the operating system's Don't-Fragment (DF) socket bit asserted.

We implemented HPQV as a complete userspace prototype in Python 3.11 with low-level C bindings to \texttt{liboqs} and \texttt{libopus}, accompanied by an interactive graphical demonstration interface capable of dynamic in-call loss injection. The protocol was evaluated through extensive empirical packet capture under Linux \texttt{tc netem} kernel impairment, an isolated Opus forward error correction benchmark, and a secondary corroborating discrete-event simulation in NS-3. Empirical packet traces verify that HPQV generates exactly zero IP fragments across all evaluated loss tiers (0\% to 30\%). Under 20\% packet loss, HPQV achieves a 99\% handshake completion rate, maintaining dependable call setup by absorbing packet drop through TCP retransmissions rather than session collapse. Furthermore, Opus in-band FEC improves audio reconstruction quality by 24 dB to 30 dB in segmental SNR compared to packet loss concealment alone. These findings demonstrate that decoupling cryptographic negotiation from datagram voice streaming provides a practical, robust, and standards-aligned pathway for deploying post-quantum security in real-time communications without sacrificing network reliability.
}
